% GENERATED FILE. Edit canonical lessons, then run npm run generate:books.
% canonical-source-hash: fnv1a64:db6885827fd0df7a
% canonical-lessons: HI-C351-adhikar, HI-C351-kartavya, HI-C351-palankarna, HI-C351-ucit, HI-C351-nyay

\chapter{Rights and Duties}
\label{ch:hi-rights-and-duties}
\hlchaptermodality{351}

\begin{chapteropening}
\textbf{By the end of this chapter:} \emph{I can talk about rights and duties.}

Complete the last lesson of chapter 351: I can talk about rights and duties.
\end{chapteropening}

\section[\texorpdfstring{\hi{अधिकार} (adhikār)}{अधिकार (adhikār)}]{\hi{अधिकार} (adhikār) — a right}

\label{lesson:HI-C351-adhikar}



\begin{quote}
Before the new one: say the Hindi for explicit, then the Hindi for to list.
\end{quote}

\subsection*{You'll want to know: \hi{अधिकार}}
\textbf{\hi{अधिकार}} — \emph{adhikār} — \textquotedblleft{}a right\textquotedblright{}.

\begin{quote}
\hi{हर} \hi{बच्चे} \hi{को} \hi{पढ़ने} \hi{का} \hi{अधिकार} \hi{है।}
\end{quote}

\begin{quote}
\emph{har bacce ko paṛhne kā adhikār hai.}
\end{quote}

\begin{quote}
Every child has the right to study.
\end{quote}

Masculine.

\begin{cousinweb}
\textbf{\hi{अधिकार}} (\emph{adhikār}) is a tatsama, borrowed from Sanskrit \emph{adhikāra}, \textquotedblleft{}authority, jurisdiction\textquotedblright{} (Monier-Williams). Hindi has narrowed it to a right that someone holds.
\end{cousinweb}

\begin{grammarlens}[title={What you've built}]
One more word for giving reasons and opinions.
\end{grammarlens}

\subsection*{Your turn}
\begin{itemize}
\raggedright
  \item \emph{Say it:} \emph{adhikār}
  \item \emph{Say it:} \emph{adhikār}, once more
  \item \emph{Say it:} \emph{har bacce ko paṛhne kā adhikār hai.} — Every child has the right to study.
  \item \emph{Recall:} say the Hindi for explicit, then the Hindi for to list, then say \emph{adhikār} again
\end{itemize}

\begin{tcolorbox}[breakable,colback=teal!4,colframe=teal!35!black,title={Before you move on}]
What does \hi{अधिकार} mean? (\textquotedblleft{}A right\textquotedblright{}.) Say it once more.
\end{tcolorbox}
\section[\texorpdfstring{\hi{कर्तव्य} (kartavya)}{कर्तव्य (kartavya)}]{\hi{कर्तव्य} (kartavya) — a duty}

\label{lesson:HI-C351-kartavya}



\begin{quote}
Before the new one: say the Hindi for to list, then the Hindi for a right.
\end{quote}

\subsection*{You'll want to know: \hi{कर्तव्य}}
\textbf{\hi{कर्तव्य}} — \emph{kartavya} — \textquotedblleft{}a duty\textquotedblright{}.

\begin{quote}
\hi{नियम} \hi{मानना} \hi{हमारा} \hi{कर्तव्य} \hi{है।}
\end{quote}

\begin{quote}
\emph{niyam mānnā hamārā kartavya hai.}
\end{quote}

\begin{quote}
Following the rules is our duty.
\end{quote}

Masculine.

\begin{cousinweb}
\textbf{\hi{कर्तव्य}} (\emph{kartavya}) is a tatsama, borrowed from Sanskrit \emph{kartavya}, \textquotedblleft{}what ought to be done\textquotedblright{} (Monier-Williams).
\end{cousinweb}

\begin{grammarlens}[title={What you've built}]
One more word for giving reasons and opinions.
\end{grammarlens}

\subsection*{Your turn}
\begin{itemize}
\raggedright
  \item \emph{Say it:} \emph{kartavya}
  \item \emph{Say it:} \emph{kartavya}, once more
  \item \emph{Say it:} \emph{niyam mānnā hamārā kartavya hai.} — Following the rules is our duty.
  \item \emph{Recall:} say the Hindi for to list, then the Hindi for a right, then say \emph{kartavya} again
\end{itemize}

\begin{tcolorbox}[breakable,colback=teal!4,colframe=teal!35!black,title={Before you move on}]
What does \hi{कर्तव्य} mean? (\textquotedblleft{}A duty\textquotedblright{}.) Say it once more.
\end{tcolorbox}
\section[\texorpdfstring{\hi{पालन} \hi{करना} (pālan karnā)}{पालन करना (pālan karnā)}]{\hi{पालन} \hi{करना} (pālan karnā) — to abide by, to follow (a rule)}

\label{lesson:HI-C351-palankarna}



\begin{quote}
Before the new one: say the Hindi for a right, then the Hindi for a duty.
\end{quote}

\subsection*{You'll want to know: \hi{पालन} \hi{करना}}
\textbf{\hi{पालन} \hi{करना}} — \emph{pālan karnā} — \textquotedblleft{}to abide by, to follow (a rule)\textquotedblright{}.

\begin{quote}
\hi{हमें} \hi{नियमों} \hi{का} \hi{पालन} \hi{करना} \hi{चाहिए।}
\end{quote}

\begin{quote}
\emph{hamẽ niyamoṁ kā pālan karnā cāhiye.}
\end{quote}

\begin{quote}
We should follow the rules.
\end{quote}

Past, with \hi{ने} (he or she …): \hi{उसने} \hi{नियम} \hi{का} \hi{पालन} \hi{किया}, \emph{usne niyam kā pālan kiyā}.

\begin{cousinweb}
\textbf{\hi{पालन} \hi{करना}} (\emph{pālan karnā}) joins \hi{पालन} to \hi{करना}, \textquotedblleft{}to do\textquotedblright{}. \hi{पालन} is a tatsama, borrowed from Sanskrit \emph{pālana}, \textquotedblleft{}guarding, protecting\textquotedblright{} (Monier-Williams). Hindi uses it for keeping to a rule.
\end{cousinweb}

\begin{grammarlens}[title={What you've built}]
One more word for giving reasons and opinions.
\end{grammarlens}

\subsection*{Your turn}
\begin{itemize}
\raggedright
  \item \emph{Say it:} \emph{pālan karnā}
  \item \emph{Say it:} \emph{pālan karnā}, once more
  \item \emph{Say it:} \emph{hamẽ niyamoṁ kā pālan karnā cāhiye.} — We should follow the rules.
  \item \emph{Recall:} say the Hindi for a right, then the Hindi for a duty, then say \emph{pālan karnā} again
\end{itemize}

\begin{tcolorbox}[breakable,colback=teal!4,colframe=teal!35!black,title={Before you move on}]
What does \hi{पालन} \hi{करना} mean? (\textquotedblleft{}To abide by, to follow (a rule)\textquotedblright{}.) Say it once more.
\end{tcolorbox}
\section[\texorpdfstring{\hi{उचित} (ucit)}{उचित (ucit)}]{\hi{उचित} (ucit) — proper, fitting}

\label{lesson:HI-C351-ucit}



\begin{quote}
Before the new one: say the Hindi for a duty, then the Hindi for to abide by.
\end{quote}

\subsection*{You'll want to know: \hi{उचित}}
\textbf{\hi{उचित}} — \emph{ucit} — \textquotedblleft{}proper, fitting\textquotedblright{}.

\begin{quote}
\hi{यह} \hi{सवाल} \hi{पूछना} \hi{उचित} \hi{नहीं} \hi{है।}
\end{quote}

\begin{quote}
\emph{yah savāl pūchnā ucit nahīṃ hai.}
\end{quote}

\begin{quote}
It isn't proper to ask this question.
\end{quote}

\begin{cousinweb}
\textbf{\hi{उचित}} (\emph{ucit}) is a tatsama, borrowed from Sanskrit \emph{ucita}, \textquotedblleft{}suitable, proper\textquotedblright{} (Monier-Williams).
\end{cousinweb}

\begin{grammarlens}[title={What you've built}]
One more word for giving reasons and opinions.
\end{grammarlens}

\subsection*{Your turn}
\begin{itemize}
\raggedright
  \item \emph{Say it:} \emph{ucit}
  \item \emph{Say it:} \emph{ucit}, once more
  \item \emph{Say it:} \emph{yah savāl pūchnā ucit nahīṃ hai.} — It isn't proper to ask this question.
  \item \emph{Recall:} say the Hindi for a duty, then the Hindi for to abide by, then say \emph{ucit} again
\end{itemize}

\begin{tcolorbox}[breakable,colback=teal!4,colframe=teal!35!black,title={Before you move on}]
What does \hi{उचित} mean? (\textquotedblleft{}Proper, fitting\textquotedblright{}.) Say it once more.
\end{tcolorbox}
\section[\texorpdfstring{\hi{न्याय} (nyāy)}{न्याय (nyāy)}]{\hi{न्याय} (nyāy) — justice}

\label{lesson:HI-C351-nyay}



\begin{quote}
Before the new one: say the Hindi for to abide by, then the Hindi for proper.
\end{quote}

\subsection*{You'll want to know: \hi{न्याय}}
\textbf{\hi{न्याय}} — \emph{nyāy} — \textquotedblleft{}justice\textquotedblright{}.

\begin{quote}
\hi{ग़रीबों} \hi{को} \hi{भी} \hi{न्याय} \hi{मिलना} \hi{चाहिए।}
\end{quote}

\begin{quote}
\emph{ġarīboṁ ko bhī nyāy milnā cāhiye.}
\end{quote}

\begin{quote}
The poor should get justice too.
\end{quote}

Masculine.

\begin{cousinweb}
\textbf{\hi{न्याय}} (\emph{nyāy}) is a tatsama, borrowed from Sanskrit \emph{nyāya}, \textquotedblleft{}a rule; a judgement\textquotedblright{} (Monier-Williams).
\end{cousinweb}

\begin{grammarlens}[title={What you've built}]
One more word for giving reasons and opinions.
\end{grammarlens}

\subsection*{Your turn}
\begin{itemize}
\raggedright
  \item \emph{Say it:} \emph{nyāy}
  \item \emph{Say it:} \emph{nyāy}, once more
  \item \emph{Say it:} \emph{ġarīboṁ ko bhī nyāy milnā cāhiye.} — The poor should get justice too.
  \item \emph{Recall:} say the Hindi for to abide by, then the Hindi for proper, then say \emph{nyāy} again
\end{itemize}

\begin{tcolorbox}[breakable,colback=teal!4,colframe=teal!35!black,title={Before you move on}]
What does \hi{न्याय} mean? (\textquotedblleft{}Justice\textquotedblright{}.) Say it once more.
\end{tcolorbox}
